\section{Discussion}
\label{sec:discussion}

\paragraph{Log loss and F1 answer different questions.}
File~10 lowers log loss on later years while its validation F1 goes down in
Phases~I and~II. That is not a failed model. F1 cuts every probability at
0.5. A model that pulls uncertain trials toward the recent success rate can
cross that cutoff the wrong way and still assign the better probability.
The contest scores the probability. We treated F1 as a check, not as the
thing to maximize. Had F1 been the target, the right next step would have
been a separate cutoff for each phase.

\paragraph{The test set is a later cohort.}
Validation trials are drawn from the same years as the training trials. Test
trials start later, and success becomes more common over time. A model fit to
the average training year is too pessimistic about the test file. Recency
weights are the direct answer: a recent trial counts more, with a decay
chosen on time-ordered holdouts. The neighbor change does the same job from
the other side. A similar past trial is more relevant if it is also recent,
and a Phase~I trial should be compared with Phase~I. File~10's validation
log loss is worse than the plain recency stack, and its Kaggle loss is
better. We trusted the check that looks like the test set.

\paragraph{Why 30\% of the old ensemble, and then stop.}
The older ensemble is weaker on later years and stronger when the years are
mixed. Blending it with file~10 is useful only while the two models miss
different trials. The first 20\% was worth 0.00135 on Kaggle. The next 10\%
was worth 0.00009. File~15 added a recent tree and a phase trend, improved
the local averages, and lost 0.00027 on Kaggle. After that, a better
historical average was not a reason to upload. The last prepared file had
small local gains, intervals that included zero, and a slightly worse grouped
score. We left the last submission slot unused.

\paragraph{Where the errors stay.}
Failed Phase~III trials are hard to recognize. They often ran to the end and
simply missed the endpoint, so the eligibility text looks like the text of a
success, and the high Phase~III success rate pulls the probability up.
Recruitment and funding failures are more administrative, and the early
ensemble catches a few more of them. The Phase~III efficacy row moves from
8 of 9 to 6 of 9 in Table~\ref{tab:p2err}, then back to 9 of 9 under the
neighbor model. All nine of those tags are Part~1 predictions. Safety,
Funding, and Other are small in every phase. Only Phase~II Recruitment and
the three no-reason groups are large, and membership in the no-reason group
is a lookup rather than a prediction. The reason table told us to look past
eligibility text. It was not a safe place to tune a cutoff.

\paragraph{What we would not do again.}
We would not ship the model with the best random validation loss. File~03
did that and scored 0.59562. We would not add fingerprints or a raw year
slope just because the feature is available. Year fitted on 2000--2013 has
to be believed out to 2022, and drug structure is missing more often in the
test file than in training. We would not treat a 0.0001 local gain as a new
submission once the leaderboard had already contradicted that kind of gain.
